\section{Pointwise Kan extensions with displayed comma colimits}\label{sec:kan}

Kan extensions organize many of the preceding universal properties. Their
existence needs care: a formula involving a colimit is not an existence
proof for that colimit. We state the exact constructive hypothesis.

Fix presented functors $K:C\to D$ and $F:C\to E$. For every object $d$
of $D$, an object of $(K\downarrow d)$ is a pair $(c,u:Kc\to d)$.
An arrow from $(c,u)$ to $(c',u')$ is $a:c\to c'$ with $u'Ka=u$.
There is a projection to $C$ and hence a diagram with value $Fc$.
Assume a colimit $L_d$ of this diagram is supplied with injections
\[
 \iota^d_{c,u}:Fc\longrightarrow L_d
\]
and its factorization and uniqueness laws. These data, and the mediator
operation on the cocones used below, must be uniform fixed PR expressions
in their numerical parameters. In particular, they respect equality of
arrow labels $u$. The cocone law is
\[
       \iota^d_{c',u'}\,Fa=\iota^d_{c,u}
                    \quad\hbox{if }u'Ka=u.                 \tag{K1}
\]
For an infinite indexing presentation, this is a schematic colimit
interface on supplied PR families; no set of all cocones is asserted.
Precisely, each fixed PR family of cocone components used in the proof,
with its displayed numerical parameters, is accompanied by a fixed PR
mediator term uniform in those parameters and by the factorization and
uniqueness proofs. The interface does not evaluate arbitrary cocone
program codes. In the finite-table case, cocones really can be supplied
as tables, and the mediator is one uniform operation on those tables.

The required equality in labels also follows from (K1): if $u\sim u'$
with the same endpoints, $\id_c$ is an arrow $(c,u)\to(c,u')$ in
the comma category, since $u'K(\id_c)\sim u$. Equation (K1),
functoriality of $F$, and the unit law therefore give
$\iota^d_{c,u}\sim\iota^d_{c,u'}$.

\begin{theorem}[Pointwise left Kan extension]\label{thm:kan}
Under these hypotheses, the assignments $d\mapsto L_d$ form a
presented functor $L:D\to E$, and there is a transformation
$\eta:F\Rightarrow LK$ such that, for every supplied presented
$G:D\to E$, composition with $\eta$ gives reciprocal constructions
\[
 \operatorname{Nat}(L,G)\ \longleftrightarrow\
 \operatorname{Nat}(F,GK).
\]
Thus $(L,\eta)$ is a left Kan extension in the presentation-wise sense.
\end{theorem}
\begin{proof}
For $v:d\to d'$, the arrows $\iota^{d'}_{c,vu}$ are a cocone over
$(K\downarrow d)$: (K1) applies because $(vu')Ka=vu$.
Let $L(v):L_d\to L_{d'}$ be its unique mediator, so
\[
                    L(v)\iota^d_{c,u}=\iota^{d'}_{c,vu}.
                                                               \tag{K2}
\]
For $v=\id$, the identity of $L_d$ has these same composites, so
uniqueness proves $L(\id)=\id$. For composable $v,w$, both $L(wv)$
and $L(w)L(v)$ have composite $\iota^{d''}_{c,wvu}$ with every
injection. Hence they are equal. The same argument shows that equal
$v$ give equal $L(v)$, so this is a functor respecting hom equality.

Define $\eta_c=\iota^{Kc}_{c,\id_{Kc}}$. If $a:c\to c'$, then
\[
 L(Ka)\eta_c=\iota^{Kc'}_{c,Ka}
            =\iota^{Kc'}_{c',\id}\,Fa=\eta_{c'}Fa,
\]
where the middle equality is (K1) for the arrow $a$ in the comma
category. Thus $\eta$ is natural.

Given $\alpha:F\Rightarrow GK$, define a cocone to $Gd$ by
\[
                   \lambda^d_{c,u}=G(u)\alpha_c.
\]
If $u'Ka=u$, its compatibility is
$G(u')\alpha_{c'}Fa=G(u')G(Ka)\alpha_c=G(u)\alpha_c$.
Let $\bar\alpha_d:L_d\to Gd$ be its mediator. To prove naturality,
compare the composites of $G(v)\bar\alpha_d$ and
$\bar\alpha_{d'}L(v)$ with $\iota^d_{c,u}$; both are
$G(vu)\alpha_c$. Uniqueness gives the naturality equation.
At the injection labelled by the identity,
\[
                  \bar\alpha_{Kc}\eta_c=\alpha_c.           \tag{K3}
\]

Conversely, suppose $\beta:L\Rightarrow G$ and
$\beta_{Kc}\eta_c=\alpha_c$. Formula (K2) gives
$\iota^d_{c,u}=L(u)\eta_c$. Therefore
\[
 \beta_d\iota^d_{c,u}
   =\beta_dL(u)\eta_c
   =G(u)\beta_{Kc}\eta_c
   =G(u)\alpha_c.
\]
Colimit uniqueness forces $\beta_d=\bar\alpha_d$ for every $d$.
This proves both inverse laws. The constructions respect transformation
equality by congruence and the same uniqueness argument.
All formulas have used fixed PR operations and finitely many logical
steps; specialization therefore yields iPRA derivations.
\end{proof}

The right Kan theorem is obtained by reversing arrows and using displayed
comma limits over $(d\downarrow K)$. If $C,D$ are finite table categories
and $F$ has finite-set values, their comma categories are finite tables.
The algorithms of Section~\ref{sec:finite} supply all required data,
so both Kan extensions exist uniformly in iPRA. For $C$ empty, the
left Kan extension has value $[0]$ and the right has value $[1]$.
The same recipes work for infinite presentations when the specific
comma (co)limits and PR mediators are supplied; their existence has not
been inferred from smallness alone.

\subsection{Universal arrows construct adjoints}
For completeness, a useful existence criterion needs much less than an
unrestricted adjoint functor theorem. Given $R:D\to C$, suppose for
each $c$ there is a supplied initial object $(Lc,\eta_c:c\to R(Lc))$
of $(c\downarrow R)$, with a uniform PR factor operation. For
$f:c\to c'$, define $Lf$ as the unique arrow satisfying
\[
                       R(Lf)\eta_c=\eta_{c'}f.
\]
The identity and composition laws follow by comparing these composites
and using initiality. The formula is also naturality of $\eta$.
For $g:c\to Rd$ the factor operation gives the unique
$\Psi(g):Lc\to d$ with $R(\Psi(g))\eta_c=g$. Its inverse is
$h\mapsto Rh\,\eta_c$. Substitution and uniqueness show naturality
in both variables, hence $L\dashv R$ by the adjunction theorem.
This is a fully constructive initial-object criterion. Merely saying
that each comma category has some initial object supplies no uniform
PR choice of the object codes or their factors; that distinction is
precisely what the hypothesis records.
